% Document d'architecture du projet GDIZ Maintenance prédictive. Compilation : cd docs && latexmk -pdf architecture.tex
\documentclass[11pt,a4paper]{article}

\usepackage[utf8]{inputenc}
\usepackage[T1]{fontenc}
\usepackage[french]{babel}
\usepackage{lmodern}
\usepackage[default]{sourcesanspro}
\usepackage{microtype}
\usepackage[a4paper,margin=2.3cm,top=2.6cm,bottom=2.6cm,headheight=14pt]{geometry}
\usepackage{xcolor}
\usepackage{tikz}
\usetikzlibrary{arrows.meta,positioning,calc}
\usepackage{booktabs}
\usepackage{array}
\usepackage{xltabular}
\usepackage{colortbl}
\usepackage{enumitem}
\usepackage{titlesec}
\usepackage{needspace}
\usepackage{fancyhdr}
\usepackage{lastpage}
\usepackage{caption}
\usepackage{float}
\usepackage{hyperref}

% ---------------------------------------------------------------------------------------------------------------------
% Charte : une couleur par couche du système, un bleu pour la structure du document
\definecolor{bleu}{HTML}{0A5CB8}
\definecolor{encre}{HTML}{1D1D1F}
\definecolor{gris}{HTML}{6E6E73}
\definecolor{filet}{HTML}{D8D8DE}
\definecolor{donnees}{HTML}{E8F1FC}
\definecolor{modele}{HTML}{FFF1DF}
\definecolor{service}{HTML}{E3F6E8}
\definecolor{trait}{HTML}{424245}
\definecolor{alerte}{HTML}{FF9500}
\definecolor{panne}{HTML}{E5372B}

% Résultats du modèle déployé (src.train, test 2025)
\newcommand{\recallval}{82\,\%}
\newcommand{\recalltest}{81\,\%}
\newcommand{\pannestest}{283}
\newcommand{\preavis}{114~h}
\newcommand{\utiles}{23\,\%}
\newcommand{\fa}{0,18}
\newcommand{\gainnet}{8,1~M~FCFA}

\hypersetup{colorlinks=true,linkcolor=bleu,urlcolor=bleu,pdftitle={GDIZ Maintenance prédictive : document d'architecture},pdfauthor={Bill Yehouenou}}
\setlist{itemsep=3pt,topsep=4pt,leftmargin=1.4em}
\setlength{\parskip}{5pt}
\setlength{\parindent}{0pt}
\renewcommand{\arraystretch}{1.25}
\color{encre}

% Titres
\titleformat{\section}{\Large\bfseries\color{bleu}}{\thesection}{0.8em}{}[{\color{filet}\titlerule[0.6pt]}]
\titleformat{\subsection}{\large\bfseries\color{encre}}{\thesubsection}{0.8em}{}
\titlespacing*{\section}{0pt}{18pt}{10pt}
\titlespacing*{\subsection}{0pt}{12pt}{4pt}

% En-têtes et pieds de page
\pagestyle{fancy}
\fancyhf{}
\fancyhead[L]{\small\color{gris}GDIZ Maintenance prédictive}
\fancyfoot[C]{\small\color{gris}\thepage\ / \pageref*{LastPage}}
\renewcommand{\headrulewidth}{0.4pt}
\renewcommand{\headrule}{\color{filet}\hrule width\headwidth height\headrulewidth}

% Légendes et tableaux
\captionsetup{font={small,color=gris},labelfont={bf,color=gris},labelsep=period,skip=6pt}
\newcolumntype{L}[1]{>{\raggedright\arraybackslash}p{#1}}
\newcommand{\entete}{}
\newcommand{\var}[1]{\texttt{#1}}


% Schémas
\tikzset{
  fleche/.style={-{Stealth[length=2.2mm]},draw=trait,line width=0.7pt},
  legende/.style={font=\scriptsize\itshape,text=gris},
}

\begin{document}

% ---------------------------------------------------------------------------------------------------------------------
% Page de garde
\begin{titlepage}
\vspace*{5.2cm}
{\color{bleu}\rule{1.6cm}{1.1mm}}\par\vspace{0.9cm}
{\fontsize{36}{42}\selectfont\bfseries\color{encre}GDIZ Maintenance\par}
{\fontsize{36}{42}\selectfont\bfseries\color{bleu}prédictive\par}
\vfill
{\color{filet}\rule{\linewidth}{0.5pt}}\par\smallskip
{\small\textcolor{gris}{Auteur}\hspace{1.2em}Bill Yehouenou\par}
\end{titlepage}

\tableofcontents
\newpage

% ---------------------------------------------------------------------------------------------------------------------
\section{Vue d'ensemble}

Le système tient en quatre couches : les \emph{données} (simulées à partir de la météo réelle), le \emph{modèle}
(entraîné puis enregistré dans MLflow), le \emph{service} (API et tableau de bord) et le \emph{suivi} (performance
réelle et dérive). Tout passe par une base DuckDB embarquée et un registre MLflow : il n'y a pas d'autre état partagé.


\begin{table}[H]
\centering\small
\begin{tabular}{@{}L{3.6cm}L{12cm}@{}}
\toprule
\entete \textbf{Dossier} & \textbf{Rôle} \\
\midrule
\var{simulation/} & Météo réelle, physique des machines, défaillances des capteurs, écriture dans DuckDB \\
\var{src/donnees.py} & Chargement, features glissantes, cible à 120~h, découpage dans le temps \\
\var{src/evaluation.py} & Métrique événementielle (pannes annoncées, fausses alertes, préavis) et bilan en FCFA \\
\var{src/train.py} & Entraînement, calibration, choix du seuil, enregistrement et promotion dans MLflow \\
\var{src/modele.py} & Modèle en service, cause probable (SHAP), scoring des 60 derniers jours \\
\var{src/metier.py} & Règles et coûts de l'atelier, indicateurs servis au tableau de bord \\
\var{src/drift.py} & Performance réelle récente et dérive des capteurs \\
\var{app/}, \var{web/} & API FastAPI et schémas ; tableau de bord HTML, CSS et JavaScript sans framework \\
\var{tests/}, \var{notebooks/} & Tests automatisés ; exploration des données et comparaison des modèles \\
\bottomrule
\end{tabular}
\caption{Organisation du code.}
\end{table}

% ---------------------------------------------------------------------------------------------------------------------
\section{Données}

Aucune donnée réelle de capteurs n'est disponible pour l'atelier de Glo-Djigbé : les données sont donc simulées, mais
ancrées dans le réel. Le point de départ est la météo horaire observée sur le site en 2024 et 2025, téléchargée depuis
Open-Meteo : température, humidité et pluie issues de la réanalyse ERA5, poussière désertique issue du service
atmosphérique CAMS, qui capte les épisodes d'Harmattan de décembre à mars. Cette météo est ramenée au pas de 30 minutes,
puis sert de décor commun aux 50 machines de l'atelier, réparties en trois gammes aux comportements différents : des
machines rustiques plus sensibles à la poussière, des machines standard, et des machines de précision protégées par
un onduleur mais plus sensibles aux variations de tension.

De l'extérieur, on passe à l'intérieur de l'atelier comme le ferait un bâtiment réel. La toiture en tôle transmet la
chaleur avec environ trois heures de retard, les machines en marche ajoutent leur propre chaleur et une ventilation
se déclenche au-delà de 38~°C. L'air réchauffé voit son humidité relative baisser, la poussière extérieure s'infiltre
progressivement et s'ajoute aux fibres textiles soulevées par la production. Le rythme de l'usine suit trois équipes
en 3$\times$8, avec une relève moins productive et un arrêt complet le week-end. Le réseau électrique de la SBEE, enfin,
connaît des délestages, plus fréquents aux heures de pointe et d'une heure environ, et une tension qui fléchit quand
la charge augmente.

Chaque machine porte ensuite un état caché qui évolue pas à pas, comme sur une vraie ligne de production. L'outil
s'use plus vite dans la poussière et sous une tension instable, la machine s'encrasse et dissipe moins bien sa chaleur,
les roulements s'usent avec la vitesse. Cet état se lit indirectement dans les capteurs : la vibration augmente avec
l'usure, la température avec l'encrassement. Cinq risques de panne se font concurrence, sur le modèle du jeu de
référence AI4I 2020 : la rupture de l'outil à l'approche d'un seuil propre à chaque machine, la surchauffe soutenue,
la surcharge mécanique, le choc électrique au retour du courant et une part d'aléatoire. Chaque panne immobilise la
machine pour une réparation de durée réaliste et remplace le composant en cause ; des entretiens préventifs, plus
probables sur une machine usée, remettent aussi les compteurs à zéro. Sur deux ans, l'atelier connaît 713 pannes.

Enfin, le modèle ne voit que ce qu'une usine verrait réellement. L'état caché n'est jamais exporté : seuls les
capteurs, les conditions de l'atelier et les compteurs de maintenance (GMAO) le sont. Les capteurs eux-mêmes sont
imparfaits, avec des paquets perdus (environ 1\,\% de valeurs manquantes), des valeurs figées et des dérives lentes.
L'ensemble, reproductible grâce à une graine aléatoire fixe, forme 1,75 million de relevés écrits dans une base
DuckDB. Ces mécanismes sont plausibles et calibrés avec soin, mais ils restent une simulation : les résultats du modèle
devront être confirmés sur de vrais capteurs.

% ---------------------------------------------------------------------------------------------------------------------
\section{Du relevé à la prédiction}

\subsection{Variables du modèle}

Le modèle voit 35 variables. Les statistiques glissantes sont calculées sur les \emph{seuls relevés en marche} : le
week-end, les capteurs tombent à zéro et feraient croire, le lundi matin, à une machine plus calme qu'elle ne l'est.
Une même fonction calcule les variables pour l'entraînement, le scoring et l'API : une prédiction n'a besoin que des
24 dernières heures de la machine.

\begin{table}[H]
\centering\small
\begin{tabular}{@{}L{4.2cm}L{9.8cm}r@{}}
\toprule
\entete \textbf{Groupe} & \textbf{Variables} & \textbf{Nombre} \\
\midrule
Gamme & \var{type\_machine} & 1 \\
Mesures du relevé & capteurs de la machine, conditions de l'atelier et réseau électrique & 11 \\
Compteurs GMAO & heures depuis le changement d'outil, depuis la dernière intervention & 2 \\
Écart thermique & température du process moins température ambiante & 1 \\
Statistiques glissantes & pour la vibration, la température du process, l'écart thermique, le couple et la poussière :
moyenne, maximum et écart-type sur 24~h, hausse sur 4~h & 20 \\
\midrule
\textbf{Total} & & \textbf{35} \\
\bottomrule
\end{tabular}
\caption{Variables d'entrée du modèle.}
\end{table}

\subsection{Cible}

Un relevé est positif si une panne prévisible (TWF, OSF ou HDF) survient sur la machine dans les 120 heures qui
suivent. Le modèle n'apprend et ne note que les relevés en marche : une machine arrêtée ne déclenche jamais d'alerte.

\subsection{Découpage dans le temps}

\begin{figure}[H]
\centering
\begin{tikzpicture}[x=0.62cm,y=0.6cm]
  \foreach \m/\l in {0/janv. 2024,3/avr.,6/juil.,9/oct.,12/janv. 2025,24/janv. 2026} {
    \draw[gris] (\m,-0.15) -- (\m,0.15);
    \node[font=\scriptsize,text=gris,below] at (\m,-0.2) {\l};
  }
  \draw[gris] (0,0) -- (24,0);
  \fill[donnees,draw=trait] (0,0.5) rectangle (3,1.3);
  \node[font=\scriptsize] at (1.5,0.9) {train};
  \fill[modele,draw=trait] (3,0.5) rectangle (6,1.3);
  \node[font=\scriptsize] at (4.5,0.9) {T2};
  \fill[modele,draw=trait] (6,0.5) rectangle (9,1.3);
  \node[font=\scriptsize] at (7.5,0.9) {T3};
  \fill[modele,draw=trait] (9,0.5) rectangle (12,1.3);
  \node[font=\scriptsize] at (10.5,0.9) {T4};
  \fill[service,draw=trait] (12,0.5) rectangle (24,1.3);
  \node[font=\scriptsize] at (18,0.9) {test 2025 : rapport uniquement};
  \draw[fleche,bleu!60] (0,1.9) -- node[above,font=\scriptsize]{modèle du pli T2} (3,1.9);
  \draw[fleche,bleu!60] (0,2.7) -- node[above,font=\scriptsize]{modèle du pli T3} (6,2.7);
  \draw[fleche,bleu!60] (0,3.5) -- node[above,font=\scriptsize]{modèle livré (calibré sur T4)} (9,3.5);
  \node[legende,align=left,anchor=west] at (12.3,2.7) {Tampon de 120~h avant chaque coupure :\\la cible d'un relevé ne doit pas\\dépendre de la période suivante.};
\end{tikzpicture}
\caption{Chaque trimestre de validation est noté par un modèle entraîné sur ce qui le précède, pour comparer les
modèles. Le modèle livré est calibré et seuillé sur T4 ; 2025 ne sert qu'au rapport.}
\end{figure}

Une machine sur cinq n'est jamais vue à l'entraînement : elle mesure la généralisation à une machine neuve.

\subsection{Modèle}

\begin{itemize}
  \item \textbf{LightGBM}, retenu face à une régression logistique, un arbre de décision et HistGradientBoosting
        (\var{notebooks/1\_comparaison\_modeles.ipynb}) : les quatre se tiennent à moins de 1,5 écart-type, LightGBM est
        le plus rapide et SHAP l'explique exactement. Réglage volontairement contraint : 80 arbres de 7 feuilles, au
        moins 1\,000 relevés par feuille ; un modèle plus libre apprend le train par cœur.
  \item \textbf{Calibration sigmoïde} sur la validation, appliquée aux log-odds : une probabilité de 10\,\% signifie
        qu'environ une machine sur dix tombe en panne dans les 5 jours.
  \item \textbf{Seuil d'alerte} : celui qui annonce le plus de pannes sans dépasser 0,2 fausse alerte par machine et
        par jour, choisi sur le dernier trimestre de validation. Il est enregistré avec le modèle : les deux ne
        peuvent pas diverger.
  \item \textbf{Cause probable} : les contributions SHAP sont regroupées en quatre familles physiques (usure et
        entretien, surchauffe, surcharge mécanique, environnement et réseau).
\end{itemize}

\subsection{Mesure de la performance}

Le système note chaque machine toutes les 30 minutes, mais un technicien ne se déplace qu'une fois pour une alerte qui
dure : la mesure compte donc des \emph{épisodes} d'alerte, redécoupés toutes les 24~h pour qu'alerter en permanence
ne paie pas.

\begin{figure}[H]
\centering
\begin{tikzpicture}[x=0.075cm,y=0.6cm]
  \draw[gris] (0,0) -- (200,0);
  \fill[panne!10] (60,0) rectangle (180,2.2);
  \node[font=\scriptsize,text=panne] at (120,2.5) {fenêtre de 120~h avant la panne};
  \draw[panne,line width=1.2pt] (180,0) -- (180,2.2);
  \node[font=\scriptsize\bfseries,text=panne,right] at (181,1.4) {panne};
  \fill[gris!50] (10,0.4) rectangle (24,1.0);
  \node[font=\scriptsize,above] at (17,1.0) {fausse};
  \fill[alerte] (95,0.4) rectangle (119,1.0);
  \fill[alerte] (119.5,0.4) rectangle (143,1.0);
  \fill[alerte] (143.5,0.4) rectangle (165,1.0);
  \node[font=\scriptsize,above] at (130,1.0) {utiles (même alerte, tranches de 24~h)};
  \draw[{Stealth[length=1.6mm]}-{Stealth[length=1.6mm]},trait] (95,-0.45) -- node[below,font=\scriptsize]{préavis} (180,-0.45);
\end{tikzpicture}
\caption{Une panne est annoncée si un épisode chevauche les 120~h qui la précèdent. Toutes les tranches de cette
fenêtre lui reviennent ; le préavis part de la première. Un épisode rattaché à aucune panne est une fausse alerte.}
\end{figure}

\begin{table}[H]
\centering\small
\begin{tabular}{@{}L{9.5cm}r@{}}
\toprule
\entete \textbf{Indicateur, modèle déployé} & \textbf{Valeur} \\
\midrule
Pannes prévisibles annoncées, validation 2024 & \recallval \\
Pannes prévisibles annoncées, test 2025 (\pannestest{} pannes) & \recalltest \\
\quad dont usure de l'outil / surchauffe / surcharge mécanique & 100 / 71 / 66\,\% \\
Préavis médian & \preavis \\
Épisodes d'alerte utiles & \utiles \\
Fausses alertes par machine et par jour (budget 0,2) & \fa \\
Gain net 2025 & \gainnet \\
\bottomrule
\end{tabular}
\caption{Performance du modèle déployé.}
\end{table}

% ---------------------------------------------------------------------------------------------------------------------
\section{Registre et cycle de vie du modèle}

Chaque entraînement enregistre dans MLflow le modèle calibré, ses hyperparamètres, son seuil exact et ses métriques.
Le service ne charge jamais autre chose que l'alias \var{@champion} : un entraînement raté ne touche pas la
production.


Le scoring écrit ensuite la probabilité de chaque relevé en marche des 60 derniers jours dans la table
\var{predictions}, avec la cause probable du dernier relevé des machines en alerte.

% ---------------------------------------------------------------------------------------------------------------------
\Needspace{28\baselineskip}  % titre, sous-titre et tableau des routes ensemble
\section{Service}

\subsection{API}

\begin{table}[H]
\centering\small
\begin{tabular}{@{}lL{5.4cm}L{8cm}@{}}
\toprule
\entete \textbf{Méthode} & \textbf{Route} & \textbf{Rôle} \\
\midrule
POST & \var{/api/v1/predict} & Un relevé : probabilité de panne et alerte \\
POST & \var{/api/v1/predict/batch} & Plusieurs relevés en un appel au modèle \\
GET & \var{/api/v1/etat} & Statut actuel des 50 machines, risque et criticité \\
GET & \var{/api/v1/parc} & Indicateurs d'une période et de la précédente, coût des pannes \\
GET & \var{/api/v1/machines/\{id\}} & Fiche machine : santé, cause, contributions SHAP, capteurs \\
GET & \var{/api/v1/machines/\{id\}/risque} & Risque au pas de 2~h sur 30 jours \\
GET / POST & \var{/api/v1/info}, \var{/reload} & Version chargée ; rechargement du @champion \\
\bottomrule
\end{tabular}
\caption{Routes de l'API (documentation interactive sur \var{/docs}).}
\end{table}

Une prédiction ajoute le relevé reçu aux 24 dernières heures de sa machine, calcule les variables avec le même code
qu'à l'entraînement et répond en une trentaine de millisecondes.

\subsection{Tableau de bord}

Servi par la même application FastAPI, en HTML, CSS et JavaScript sans framework :
\begin{itemize}
  \item \textbf{Parc} : indicateurs de la période (disponibilité, MTBF, coût des pannes, heures d'arrêt) comparés à la
        période précédente, plan de l'atelier, matrice risque $\times$ criticité qui dit quoi faire (agir, planifier,
        surveiller), pannes par gamme et par cause, coût par mois ;
  \item \textbf{Machine} : santé, risque à 5 jours, cause probable expliquée en une phrase, contributions SHAP,
        historique des pannes sur 12 mois, capteurs face aux limites du parc, courbe du risque sur 30 jours.
\end{itemize}

La criticité vient de la chaîne de production : une cardeuse (4 machines sur l'étape) pèse plus qu'un métier à tisser
(24 machines). Le coût d'une panne compte les pièces, la production perdue et le technicien pendant la durée réelle
de l'arrêt, et le tissu rebuté ; une panne annoncée devient un arrêt planifié de 2~h sans rebut.

% ---------------------------------------------------------------------------------------------------------------------
\section{Suivi en production}

\var{make check-drift} répond à deux questions :
\begin{itemize}
  \item \textbf{Le modèle tient-il ?} Sur les 60 derniers jours, les pannes passées sont connues : on les compare aux
        alertes qui les précédaient, avec la même métrique qu'à l'entraînement.
  \item \textbf{Les capteurs dérivent-ils ?} Les 30 derniers jours sont comparés aux mêmes jours un an plus tôt
        (indice de stabilité de population, PSI) : la météo est saisonnière, comparer avec le mois précédent
        signalerait l'Harmattan comme une dérive.
\end{itemize}
Une dérive est un signal à examiner, pas un réentraînement automatique.

% ---------------------------------------------------------------------------------------------------------------------
\Needspace{10\baselineskip}
\section{Limites}

\begin{enumerate}
  \item Les données sont simulées : les mécanismes sont plausibles, mais les résultats restent à confirmer sur de vrais
        capteurs.
  \item Les pannes électriques et aléatoires n'ont pas de signe avant-coureur et ne sont pas apprises.
  \item Moins d'une alerte sur quatre précède une panne : le modèle n'est rentable que tant qu'une inspection coûte
        moins d'environ 5\,700~FCFA.
  \item La cause probable désigne une famille physique, pas un diagnostic du mode de panne.
\end{enumerate}

% ---------------------------------------------------------------------------------------------------------------------
\section*{Glossaire}
\addcontentsline{toc}{section}{Glossaire}

{\small
\begin{xltabular}{\linewidth}{@{}L{3.4cm}X@{}}
\toprule
\entete \textbf{Terme} & \textbf{Définition} \\
\midrule
Calibration & Ajustement des scores du modèle pour qu'ils se lisent comme des probabilités \\
Champion & Version du modèle servie en production (alias MLflow \var{@champion}) \\
Épisode d'alerte & Suite d'alertes consécutives sur une machine, redécoupée toutes les 24~h : une intervention \\
GMAO & Gestion de maintenance assistée par ordinateur : compteurs d'heures et historique des interventions \\
Harmattan & Vent sec chargé de poussière saharienne, de décembre à mars en Afrique de l'Ouest \\
MTBF, MTTR & Temps moyen de marche entre deux pannes ; temps moyen de réparation \\
PSI & Population Stability Index : écart entre deux distributions d'une même variable \\
Recall événementiel & Part des pannes précédées d'au moins une alerte dans les 120~h \\
SBEE & Société béninoise d'énergie électrique, dont les délestages coupent le courant \\
SHAP & Méthode qui attribue à chaque variable sa contribution à une prédiction \\
TWF, OSF, HDF, PWF, RNF & Modes de panne : usure de l'outil, surcharge mécanique, surchauffe, électrique, aléatoire \\
\bottomrule
\end{xltabular}}

\end{document}
